\newpage
\phantomsection
\label{prf:differentiation-on-set}

\begin{remark*}[Return]
\hyperref[thm:finite-linear-combination-rule-interval]{Return to Theorem}
\end{remark*}

\begin{theorem*}[Differentiation Algebra on Sets]
Let $I \subseteq \mathbb{R}$ be an open interval and let $f, g : I \to \mathbb{R}$ be differentiable on $I$. Let $\alpha \in \mathbb{R}$.
\begin{enumerate}
  \item \textbf{Constant Multiple Rule:} $(\alpha f)' = \alpha f'$ on $I$.
  \item \textbf{Sum Rule:} $(f + g)' = f' + g'$ on $I$.
  \item \textbf{Product Rule:} $(fg)' = f'g + fg'$ on $I$.
  \item \textbf{Quotient Rule:} If $g(x) \neq 0$ for all $x \in I$, then $(f/g)' = (f'g - fg')/g^2$ on $I$.
\end{enumerate}
\end{theorem*}

\begin{proof}
\textbf{Professional Standard Proof.}~\\
Each result follows by pointwise application of the corresponding differentiation rule at arbitrary points $c \in I$.

\textbf{Constant Multiple Rule:} Let $h := \alpha f$ and fix $c \in I$. Since $f$ is differentiable at $c$, the pointwise Constant Multiple Rule gives $h'(c) = \alpha f'(c)$. Since $c \in I$ was arbitrary, $h' = \alpha f'$ on $I$.

\textbf{Sum Rule:} Let $h := f + g$ and fix $c \in I$. Since $f$ and $g$ are differentiable at $c$, the pointwise Sum Rule gives $h'(c) = f'(c) + g'(c)$. Since $c \in I$ was arbitrary, $h' = f' + g'$ on $I$.

\textbf{Product Rule:} Let $h := fg$ and fix $c \in I$. Since $f$ and $g$ are differentiable at $c$, the pointwise Product Rule gives $h'(c) = f'(c)g(c) + f(c)g'(c)$. Since $c \in I$ was arbitrary, $h' = f'g + fg'$ on $I$.

\textbf{Quotient Rule:} Let $h := f/g$ and fix $c \in I$. Since $g(c) \neq 0$ and $f, g$ are differentiable at $c$, the pointwise Quotient Rule gives $h'(c) = (f'(c)g(c) - f(c)g'(c))/(g(c))^2$. Since $c \in I$ was arbitrary and $g(x) \neq 0$ on all of $I$, $h' = (f'g - fg')/g^2$ on $I$.
\end{proof}

\begin{proof}
\textbf{Detailed Learning Proof.}~\\

\textbf{Step 1.} Establish the universal strategy. Each global differentiation rule is obtained by applying the corresponding pointwise rule at an arbitrary point $c \in I$ and concluding the function identity holds on the entire interval $I$.

\textbf{Step 2.} Prove the Constant Multiple Rule. Let $h := \alpha f$ and fix arbitrary $c \in I$. Since $f$ is differentiable at $c$, the pointwise Constant Multiple Rule ensures $h$ is differentiable at $c$ with $h'(c) = \alpha f'(c)$. Since $c \in I$ was arbitrary, this equality holds at every point of $I$, giving $h' = \alpha f'$ on $I$.

\textbf{Step 3.} Prove the Sum Rule. Let $h := f + g$ and fix arbitrary $c \in I$. Since both $f$ and $g$ are differentiable at $c$, the pointwise Sum Rule ensures $h$ is differentiable at $c$ with $h'(c) = f'(c) + g'(c)$. Since $c \in I$ was arbitrary, this equality holds at every point of $I$, giving $h' = f' + g'$ on $I$.

\textbf{Step 4.} Prove the Product Rule. Let $h := fg$ and fix arbitrary $c \in I$. Since both $f$ and $g$ are differentiable at $c$, the pointwise Product Rule ensures $h$ is differentiable at $c$ with $h'(c) = f'(c)g(c) + f(c)g'(c)$. Since $c \in I$ was arbitrary, this equality holds at every point of $I$, giving $h' = f'g + fg'$ on $I$.

\textbf{Step 5.} Prove the Quotient Rule. Let $h := f/g$ and fix arbitrary $c \in I$. The global hypothesis ensures $g(c) \neq 0$, so the pointwise Quotient Rule applies: $h$ is differentiable at $c$ with $h'(c) = (f'(c)g(c) - f(c)g'(c))/(g(c))^2$. Since $c \in I$ was arbitrary and $g(x) \neq 0$ throughout $I$, this equality holds at every point of $I$, giving $h' = (f'g - fg')/g^2$ on $I$.
\end{proof}

\begin{remark*}[Proof structure]
Universal instantiation strategy. Each result lifts the corresponding pointwise differentiation rule to a global statement on the interval $I$ by applying the pointwise rule at an arbitrary point $c \in I$ and leveraging the arbitrariness to conclude the function identity holds throughout $I$. The Quotient Rule requires the additional global hypothesis that $g$ is nonzero on all of $I$ to ensure the pointwise condition is satisfied at every application.
\end{remark*}

\begin{remark*}[Dependencies]~\\
\begin{itemize}
  \item \hyperref[thm:constant-multiple-rule]{Constant Multiple Rule (Pointwise)}
  \item \hyperref[thm:sum-rule]{Sum Rule (Pointwise)}
  \item \hyperref[thm:product-rule]{Product Rule (Pointwise)}
  \item \hyperref[thm:quotient-rule]{Quotient Rule (Pointwise)}
\end{itemize}
\end{remark*}